% ભાગ: first-order-logic
% પ્રકરણ: beyond
% વિભાગ: second-order-logic

\documentclass[../../../include/open-logic-section]{subfiles}

\begin{document}

\olfileid{fol}{byd}{sol}

\olsection{દ્વિતીય-ક્રમ તર્કશાસ્ત્ર}

દ્વિતીય-ક્રમ તર્કશાસ્ત્રની ભાષા માત્ર વ્યક્તિઓના !!a{domain} પર જ નહિ,
પણ એ !!{domain} પરના સંબંધો પર પણ પરિમાણીકરણ કરવા દે છે. કોઈ
પ્રથમ-ક્રમ ભાષા~$\Lang{L}$ આપેલી હોય, તો દરેક~$k$ માટે $k$-સ્થાની
સંબંધો પર વ્યાપતા !!{variable}s~$R$ ઉમેરવામાં આવે છે અને એ ચલો પર
પરિમાણીકરણની છૂટ આપવામાં આવે છે. જો~$R$ કોઈ $k$-સ્થાની સંબંધ માટેનો
!!a{variable} હોય અને $t_1$, \dots,~$t_k$ સામાન્ય (પ્રથમ-ક્રમ) પદો હોય,
તો $\Atom{R}{t_1,\dots,t_k}$ આણ્વિક !!{formula} છે. બાકી !!{formula}sનો
ગણ પ્રથમ-ક્રમ તર્કશાસ્ત્રની જેમ જ વ્યાખ્યાયિત થાય છે, પરંતુ તેમાં
દ્વિતીય-ક્રમ પરિમાણીકરણ માટે વધારાની શરતો હોય છે. નોંધો કે આપણી પાસે
!!{identity} માત્ર પ્રથમ-ક્રમ પદો માટે છે: જો $R$ અને~$S$ સમાન
સ્થાનસંખ્યા~$k$ ધરાવતા સંબંધ-!!{variable}s હોય, તો $\eq[R][S]$ નીચેનું
સંક્ષેપ છે:
\[
\lforall[x_1 \dots][\lforall[x_k][(\Atom{R}{x_1, \dots, x_k} \liff
  \Atom{S}{x_1, \dots, x_k})]].
\]

દ્વિતીય-ક્રમ તર્કશાસ્ત્રના નિયમો પરિમાણક-નિયમોને નવા દ્વિતીય-ક્રમ
ચલો સુધી વિસ્તારે છે. જોકે અહીં આ ચલો~$\Lang{L}$ના !!{predicate}s અને
વધુ સામાન્ય રીતે~$\Lang{L}$ના !!{formula}s સાથે કેવી રીતે ક્રિયા કરે
છે તે સમજાવવામાં થોડી કાળજી લેવી પડે છે. ઓછામાં ઓછું તો સંબંધ-ચલોને
પદો ગણવામાં આવે છે, તેથી આપણી પાસે આ સ્વરૂપનાં અનુમાનો છે:
\[
!A(R) \Proves \lexists[R][!A(R)]
\]
પરંતુ જો~$\Lang{L}$ અચળ સંબંધ-સંકેત~$<$ ધરાવતી અંકગણિતની ભાષા હોય,
તો નીચેનું અનુમાન પણ પ્રામાણ્ય ધરાવે એવી અપેક્ષા રાખીશું:
\[
x < y \Proves \lexists[R][\Atom{R}{x,y}]
\]
અથવા કોઈ આપેલા !!{formula}~$!A$ માટે,
\[
!A(x_1, \dots, x_k) \Proves \lexists[R][\Atom{R}{x_1,\dots,x_k}]
\]
વધુ સામાન્ય રીતે, આ સ્વરૂપનાં અનુમાનોને છૂટ આપવા માગી શકીએ:
\[
\Subst{!A}{\lambd[\vec x][!B(\vec x)]}{R} \Proves
\lexists[R][!A]
\]
અહીં $\Subst{!A}{\lambd[\vec x][!B(\vec x)]}{R}$નો અર્થ~$!A$માં
$\Obj{R}{t_1,\dots,t_k}$ સ્વરૂપના દરેક આણ્વિક !!{formula}ને
$!B(t_1, \dots, t_k)$ વડે બદલવાથી મળતું પરિણામ છે. આ છેલ્લો નિયમ
\emph{ગુણધર્મ-ગણરચના પ્રરૂપ} રાખવા સમકક્ષ છે, એટલે આ સ્વરૂપનું
સ્વયંસિદ્ધિ રાખવા સમકક્ષ છે:
\[
\lexists[R][\lforall[x_1, \dots, x_k][(!A(x_1, \dots, x_k) \liff
\Atom{R}{x_1, \dots, x_k})]],
\]
દ્વિતીય-ક્રમ ભાષાના દરેક એવા !!{formula}~$!A$ માટે જેમાં~$R$ મુક્ત
ચલ નથી.
(સ્વાધ્યાય: જો~$!A$માં~$R$ને આવવાની છૂટ હોય, તો આ પ્રરૂપ અસંગત છે તે
બતાવો!)

તર્કશાસ્ત્રીઓ ``દ્વિતીય-ક્રમ તર્કશાસ્ત્રનાં સ્વયંસિદ્ધિઓ'' કહે ત્યારે
સામાન્ય રીતે દ્વિતીય-ક્રમ પરિમાણક-નિયમો વડે પ્રથમ-ક્રમ તર્કશાસ્ત્રનો
લઘુતમ વિસ્તાર અને તેની સાથે ગુણધર્મ-ગણરચના પ્રરૂપ સમજતા હોય છે. પરંતુ
આ સ્વયંસિદ્ધિઓ અને નિયમોનાં દુર્બળતર ઉપતંત્રોનો અભ્યાસ કરવો પણ ઘણી વાર
રસપ્રદ છે. દાખલા તરીકે, નોંધો કે પોતાના સંપૂર્ણ સામાન્ય સ્વરૂપમાં
ગુણધર્મ-ગણરચનાનું સ્વયંસિદ્ધિ-પ્રરૂપ \emph{અપ્રેડિકેટિવ} છે: તે
દ્વિતીય-ક્રમ પરિમાણકો ધરાવતા !!a{formula} વડે ``વ્યાખ્યાયિત'' થતા
સંબંધ~$\Atom{R}{x_1, \dots, x_k}$ના અસ્તિત્વનું વિધાન કરવા દે છે; અને
આ પરિમાણકો આવા બધા સંબંધોના ગણ પર વ્યાપે છે---જેમાં~$R$ પોતે પણ સામેલ
છે!{} વીસમી સદીની શરૂઆતની આસપાસ રસેલના વિરોધાભાસ પ્રત્યેની એક સામાન્ય
પ્રતિક્રિયા આવી વ્યાખ્યાઓને દોષ આપવાની અને ગણિતના પાયાના વિકાસમાં તેમને
ટાળવાની હતી. જો !!{formula}~$!A$માં દ્વિતીય-ક્રમ પરિમાણકોનો ઉપયોગ
નિષિદ્ધ કરીએ, તો ગુણધર્મ-ગણરચનાનું કંઈક દુર્બળ એવું
\emph{પ્રેડિકેટિવ} સ્વરૂપ મળે છે.

અર્થવિજ્ઞાનની દૃષ્ટિએ, દ્વિતીય-ક્રમ !!{structure}ને ભાષા માટેની
પ્રથમ-ક્રમ !!{structure} અને !!{domain} પરના એવા સંબંધોના ગણના જોડાણ
તરીકે વિચારી શકાય છે જેના પર દ્વિતીય-ક્રમ પરિમાણકો વ્યાપે છે (વધુ
ચોક્કસ રીતે, દરેક~$k$ માટે સ્થાનસંખ્યા~$k$ ધરાવતા સંબંધોનો એક ગણ હોય
છે). અલબત્ત, ગુણધર્મ-ગણરચના !!{derivation} તંત્રમાં સામેલ હોય, તો
ગુણધર્મ-ગણરચનાનાં સ્વયંસિદ્ધિઓ સંતોષવા માટે ``દ્વિતીય-ક્રમ ભાગ''માં
પૂરતા સંબંધો હોવાની વધારાની શરત મૂકવી પડે છે---નહિતર !!{derivation}
તંત્ર યથાર્થ નથી!{} પૂરતા સંબંધો ઉપલબ્ધ રહે તેની ખાતરી કરવાની એક સરળ
રીત દ્વિતીય-ક્રમ ભાગમાં પ્રથમ-ક્રમ ભાગ પરના \emph{બધા} સંબંધો લેવાની
છે. આવી !!a{structure}ને \emph{પૂર્ણ} કહે છે અને એક અર્થમાં તે ભાષાની
ખરી ``અભિપ્રેત !!{structure}'' છે. જો આપણે ધ્યાન પૂર્ણ !!{structure}s
સુધી મર્યાદિત કરીએ, તો તેને દ્વિતીય-ક્રમનું \emph{પૂર્ણ} અર્થવિજ્ઞાન
કહે છે. આ કિસ્સામાં !!{structure}નું નિર્દેશન માત્ર પ્રથમ-ક્રમ ભાગનું
નિર્દેશન કરવા જેટલું રહે છે, કારણ કે દ્વિતીય-ક્રમ ભાગની સામગ્રી તેમાંથી
અંતર્નિહિત રીતે નક્કી થાય છે.

સારરૂપે, દ્વિતીય-ક્રમ તર્કશાસ્ત્રની વાત કરતી વખતે થોડી સંદિગ્ધતા રહે છે.
!!{derivation} તંત્રની દૃષ્ટિએ નીચેમાંથી એક બાબત મનમાં હોઈ શકે:
\begin{enumerate}
\item કેટલાક ગુણધર્મ-ગણરચના સ્વયંસિદ્ધિઓ સાથેનું ``લઘુતમ'' દ્વિતીય-ક્રમ
  !!{derivation} તંત્ર.
\item પૂર્ણ ગુણધર્મ-ગણરચના ધરાવતું ``પ્રમાણભૂત'' દ્વિતીય-ક્રમ
  !!{derivation} તંત્ર.
\end{enumerate}
અર્થવિજ્ઞાનની દૃષ્ટિએ નીચેમાંથી એક બાબતમાં રસ હોઈ શકે:
\begin{enumerate}
\item ``દુર્બળ'' અર્થવિજ્ઞાન, જેમાં !!a{structure} પ્રથમ-ક્રમ ભાગ અને
  ગુણધર્મ-ગણરચનાનાં સ્વયંસિદ્ધિઓ સંતોષવા પૂરતો મોટો દ્વિતીય-ક્રમ ભાગ
  ધરાવે છે.
\item ``પ્રમાણભૂત'' દ્વિતીય-ક્રમ અર્થવિજ્ઞાન, જેમાં માત્ર પૂર્ણ
  !!{structure}sને વિચારવામાં આવે છે.
\end{enumerate}
તર્કશાસ્ત્રીઓ પોતાના મનમાં રહેલું !!{derivation} તંત્ર કે અર્થવિજ્ઞાન
સ્પષ્ટ ન કરે ત્યારે સામાન્ય રીતે બંને યાદીની બીજી બાબતનો ઉલ્લેખ કરતા
હોય છે. આ અર્થવિજ્ઞાનનો લાભ એ છે કે, આપણે જોઈશું તેમ, તે અનેક કુદરતી
ગાણિતિક સંરચનાઓનું એકરૂપતા સુધીનું અનન્ય વર્ણન આપે છે; સાથે સાથે
!!{derivation} તંત્ર ઘણું શક્તિશાળી અને આ અર્થવિજ્ઞાન માટે યથાર્થ છે.
તેની ખામી એ છે કે !!{derivation} તંત્ર આ અર્થવિજ્ઞાન માટે \emph{પૂર્ણ
નથી}; હકીકતમાં, અસરકારક રીતે આપેલું \emph{કોઈ} !!{derivation} તંત્ર
પૂર્ણ દ્વિતીય-ક્રમ અર્થવિજ્ઞાન માટે પૂર્ણ નથી. બીજી બાજુ, આપણે જોઈશું
કે આ !!{derivation} તંત્ર દુર્બળ અર્થવિજ્ઞાન માટે \emph{પૂર્ણ} છે; તેથી
જો કોઈ વાક્ય સાબિત કરી શકાય નહિ, તો એવી \emph{કોઈક} !!{structure}
હોય છે---જરૂરી નથી કે પૂર્ણ હોય---જેમાં તે અસત્ય છે.

દ્વિતીય-ક્રમ તર્કશાસ્ત્રની ભાષા ઘણી સમૃદ્ધ છે. એકસ્થાની સંબંધોને
!!{domain}ના ઉપગણો સાથે એકરૂપ ગણી શકાય છે, તેથી ખાસ કરીને આ ગણો પર
પરિમાણીકરણ કરી શકાય છે; દાખલા તરીકે, પ્રાકૃતિક સંખ્યાઓ માટેનું
અનુમાનપ્રવર્તન એક જ સ્વયંસિદ્ધિ વડે વ્યક્ત કરી શકાય છે:
\[
\lforall[R][((\Atom{R}{\Obj{0}} \land \lforall[x][(\Atom{R}{x} \lif
    \Atom{R}{x'})]) \lif \lforall[x][\Atom{R}{x}])].
\]
અંકગણિતની ભાષામાં $\Obj 0, \Obj \prime, +, \times$ અને~$<$ સંકેતો
હોય, તો તેમનું વર્તન વર્ણવવા નીચેની સ્વયંસિદ્ધિઓ ઉમેરી શકાય છે:
\begin{enumerate}
\item $\lforall[x][\lnot x' = \Obj 0]$
\item $\lforall[x][\lforall[y][(x' = y' \lif x = y)]]$
\item $\lforall[x][(x + \Obj 0) = x]$
\item $\lforall[x][\lforall[y][(x + y') = (x + y)']]$
\item $\lforall[x][(x \times \Obj 0) = \Obj 0]$
\item $\lforall[x][\lforall[y][(x \times y') = ((x \times y) + x)]]$
\item $\lforall[x][\lforall[y][(x < y \liff \lexists[z][y = (x + z')])]]$
\end{enumerate}
\sourcecorrection{OLBYD-002}{મૂળની
\texttt{second-order-logic.tex}, પંક્તિ~142માં બીજા સ્વયંસિદ્ધિમાં
અચાનક~$s(x)=s(y)$ લખાયું છે, જ્યારે જાહેર કરેલો ઉત્તરગામી સંકેત
$\Obj\prime$ છે અને એ જ યાદી તથા આગળની સાબિતીમાં અગ્રચિહ્ન વપરાયું છે.
અહીં સુસંગત~$x'=y'$ પુનઃસ્થાપિત કર્યું છે.}
આ સ્વયંસિદ્ધિઓ અને ઉપરના અનુમાનપ્રવર્તનના સ્વયંસિદ્ધિ સાથે મળીને, પૂર્ણ
દ્વિતીય-ક્રમ અર્થવિજ્ઞાન વાપરીએ તો, અંકગણિતના પ્રમાણભૂત નિદર્શ
!!{structure}~$\Struct{N}$નું એકરૂપતા સુધીનું અનન્ય વર્ણન આપે છે એવું
બતાવવું મુશ્કેલ નથી. આ સ્વયંસિદ્ધિઓ સત્ય હોય એવી કોઈ પણ
!!{structure}~$\Struct{M}$ આપેલી હોય, તો $\Nat$થી~$\Struct{M}$ના
!!{domain}માં વિધેય~$f$ને~$\Nat$ પરના સામાન્ય આવર્તન વડે આ રીતે
વ્યાખ્યાયિત કરો: $f(0) = \Assign{\Obj 0}{M}$ અને
$f(x+1) = \Assign{\prime}{M}(f(x))$. $\Nat$ પરના સામાન્ય
અનુમાનપ્રવર્તન અને સ્વયંસિદ્ધિઓ (1) તથા~(2)~$\Struct M$માં સત્ય છે એ
હકીકતથી $f$~!!{injective} છે તે મળે છે. $f$~!!{surjective} છે તે
જોવા માટે~$P$ને~$f$ના પરાસમાં આવતા $\Domain{M}$ના ઘટકોનો ગણ
લો. $\Struct M$ પૂર્ણ હોવાથી~$P$ દ્વિતીય-ક્રમ !!{domain}માં છે.
$f$ની રચનાથી $\Assign{\Obj 0}{M}$,~$P$માં છે અને~$P$,
$\Assign{\prime}{M}$ હેઠળ સંવૃત છે. $\Struct M$માં અનુમાનપ્રવર્તનનું
સ્વયંસિદ્ધિ સત્ય છે (ખાસ કરીને~$P$ માટે) તેથી~$P$,~$\Struct M$ના
સમગ્ર પ્રથમ-ક્રમ !!{domain} જેટલો છે. આ બતાવે છે કે~$f$~!!a{bijection}
છે. $\Nat$ પરના સામાન્ય અનુમાનપ્રવર્તનનો વારંવાર ઉપયોગ કરીને~$f$
સમરૂપતા છે તે બતાવવું વધુ મુશ્કેલ નથી.

ગણસિદ્ધાંતની ભાષામાં વિધેય માત્ર એક ખાસ પ્રકારનો સંબંધ છે; દાખલા
તરીકે, એકસ્થાની વિધેય~$f$ને
$\lforall[x][\lexists![y][R(x,y)]]$ સંતોષતા દ્વિસ્થાની સંબંધ~$R$ સાથે
એકરૂપ ગણી શકાય છે. તેથી વિધેયો પર પણ પરિમાણીકરણ કરી શકાય છે. પૂર્ણ
અર્થવિજ્ઞાન વાપરીને અનંત !!{structure}sના વર્ગને એવી
!!{structure}s~$\Struct M$ના વર્ગ તરીકે વ્યાખ્યાયિત કરી શકાય છે જેના
માટે~$\Struct M$ના !!{domain}માંથી તેના પોતાના યથાર્થ ઉપગણમાં કોઈ
એકૈક વિધેય હોય:
\[
\lexists[f][(\lforall[x][\lforall[y][(\eq[f(x)][f(y)] \lif \eq[x][y])]]
  \land \lexists[y][\lforall[x][\eq/[f(x)][y]]])].
\]
પછી આ વાક્યનો નિષેધ સાન્ત !!{structure}sના વર્ગને વ્યાખ્યાયિત કરે છે.

વળી, રેખીય ક્રમની વ્યાખ્યામાં નીચેનું ઉમેરીને સુક્રમોના વર્ગને
વ્યાખ્યાયિત કરી શકાય છે:
\[
\lforall[P][(\lexists[x][\Atom{P}{x}] \lif \lexists[x][(\Atom{P}{x}
    \land \lforall[y][(y < x \lif \lnot \Atom{P}{y})])])].
\]
``ગણ''ને ``એકસ્થાની સંબંધ'' સાથે એકરૂપ ગણીએ તો આ કહે છે કે દરેક
અરિક્ત ગણને લઘુતમ ઘટક હોય છે. બીજા દાખલા તરીકે, શિરોબિંદુઓને
અસંબદ્ધ ભાગોમાં બિનમામૂલી રીતે વહેંચી શકાય નહિ એમ કહીને આલેખોની
સંલગ્નતાનો ખ્યાલ વ્યક્ત કરી શકાય છે:
\[
\lnot \lexists[A][(\lexists[x][A(x)] \land \lexists[y][\lnot A(y)]
  \land \lforall[w][\lforall[z][((\Atom{A}{w} \land \lnot \Atom{A}{z})
      \lif \lnot \Atom{R}{w,z})]])].
\]
હજુ એક દાખલા તરીકે, જે સાન્ત !!{structure}sના !!{domain}નું કદ બેકી
હોય તેમનો વર્ગ વ્યાખ્યાયિત કરવાનો સ્વાધ્યાય કરી શકો. વધુ નોંધપાત્ર રીતે,
પરિમેય સંખ્યાઓ ધરાવતા પૂર્ણ ક્રમિત ક્ષેત્ર તરીકે વાસ્તવિક સંખ્યાઓનું
એકરૂપતા સુધીનું અનન્ય વર્ણન આપી શકાય છે.

ટૂંકમાં, દ્વિતીય-ક્રમ તર્કશાસ્ત્ર પ્રથમ-ક્રમ તર્કશાસ્ત્ર કરતાં ઘણું વધુ
અભિવ્યક્તિશીલ છે. આ સારા સમાચાર છે; હવે ખરાબ સમાચાર. પૂર્ણ દ્વિતીય-ક્રમ
અર્થવિજ્ઞાન માટે પૂર્ણ હોય એવું કોઈ અસરકારક !!{derivation} તંત્ર નથી
એનો ઉલ્લેખ આપણે કરી ચૂક્યા છીએ. સારું હોય કે ખરાબ, પ્રથમ-ક્રમ
તર્કશાસ્ત્રના સઘનતા અને લેવેનહાઇમ--સ્કોલેમ પ્રમેયો સહિતના ઘણા ગુણધર્મો
અહીં ગેરહાજર છે.

બીજી બાજુ, જો પૂર્ણ દ્વિતીય-ક્રમ અર્થવિજ્ઞાન છોડીને દુર્બળ અર્થવિજ્ઞાન
સ્વીકારવા તૈયાર હોઈએ, તો લઘુતમ દ્વિતીય-ક્રમ !!{derivation} તંત્ર આ
અર્થવિજ્ઞાન માટે પૂર્ણ છે. બીજા શબ્દોમાં, $\Proves$ને ``લઘુતમ તંત્રમાં
સાબિત કરે છે'' અને $\Entails$ને ``દુર્બળ અર્થવિજ્ઞાનમાં તાર્કિક રીતે
ફલિત કરે છે'' તરીકે વાંચીએ, તો જ્યારે પણ $\Gamma \Entails !A$, ત્યારે
$\Gamma \Proves !A$ એમ બતાવી શકીએ. !!{derivation} તંત્રમાં ચોક્કસ
ગુણધર્મ-ગણરચના સ્વયંસિદ્ધિઓ સામેલ કરવા માગીએ, તો આ સ્વયંસિદ્ધિઓ
સંતોષતી દ્વિતીય-ક્રમ સંરચનાઓ સુધી અર્થવિજ્ઞાનને મર્યાદિત કરવું પડે:
દાખલા તરીકે, જો~$\Delta$ ગુણધર્મ-ગણરચના સ્વયંસિદ્ધિઓનો કોઈ ગણ હોય
(કદાચ એ બધાં જ), તો $\Gamma \cup \Delta \Entails !A$ હોય ત્યારે
$\Gamma \cup \Delta \Proves !A$. ખાસ કરીને, આપણે વિચારતાં
ગુણધર્મ-ગણરચના સ્વયંસિદ્ધિઓ વડે~$!A$ સાબિત કરી શકાય નહિ, તો
$\lnot !A$નું એવું નિદર્શ છે જેમાં આ ગુણધર્મ-ગણરચના સ્વયંસિદ્ધિઓ છતાં
સત્ય છે.

દુર્બળ અર્થવિજ્ઞાન માટે પૂર્ણતા પ્રમેય સત્ય છે તે જોવાની સૌથી સરળ રીત
દ્વિતીય-ક્રમ તર્કશાસ્ત્રને નીચે મુજબ બહુ-પ્રકાર તર્કશાસ્ત્ર તરીકે
વિચારવાની છે. એક પ્રકારનું અર્થઘટન સામાન્ય ``પ્રથમ-ક્રમ'' !!{domain}
તરીકે થાય છે અને પછી દરેક~$k$ માટે ``સ્થાનસંખ્યા~$k$ ધરાવતા સંબંધો''નું
!!a{domain} હોય છે. ભાષામાં આંતરિક સંબંધ-સંકેતો
``$\Atom{\Obj{true}_k}{R,x_1,\dots,x_k}$'' લઈએ છીએ, જે કહેવા માટે છે
કે~$R$, $x_1$, \dots,~$x_k$ માટે સત્ય છે; અહીં~$R$ ``$k$-સ્થાની
સંબંધ'' પ્રકારનો ચલ છે અને $x_1$, \dots,~$x_k$ પ્રથમ-ક્રમ પ્રકારની
વસ્તુઓ છે.

આ એકરૂપીકરણ હેઠળ દુર્બળ દ્વિતીય-ક્રમ અર્થવિજ્ઞાન મૂળભૂત રીતે બહુ-પ્રકાર
તર્કશાસ્ત્રનું સામાન્ય અર્થવિજ્ઞાન છે; અને બહુ-પ્રકાર તર્કશાસ્ત્રને
પ્રથમ-ક્રમ તર્કશાસ્ત્રમાં સમાવી શકાય છે એવું આપણે જોઈ ચૂક્યા છીએ. તેથી
બંને દિશાના અનુવાદોને બાદ કરીએ, તો દ્વિતીય-ક્રમ તર્કશાસ્ત્રની દુર્બળ
કલ્પના ખરેખર વેશપલટો કરેલું પ્રથમ-ક્રમ તર્કશાસ્ત્ર છે, જેમાં !!{domain}
યોગ્ય સ્વયંસિદ્ધિઓ દ્વારા નિયંત્રિત ``વસ્તુઓ'' અને ``સંબંધો'' બંને ધરાવે
છે.

\end{document}
